% Section 90 --- Scope, Status, and Qualifications. Drafted after the body,
% from research/scope.md and the body sections as actually written; skeleton
% per the rosary leaf's section 90.
\sectionguard
\section{Scope, Status, and Qualifications}

\textbf{Question and coverage.} This study answers: where the ark of the
covenant went, so far as Scripture, tradition, and scholarship can say;
what its journey teaches about presence, holiness, and the propitiatory;
and how the Church came to read the Virgin Mary as the ark of the new
covenant --- in philology, Fathers, liturgy, and magisterium --- with
every parallel graded separately for attestation and reception, ending at
Apocalypse 11:19--12:1 and the Assumption. It is a source-audited study
under the repository's Mariology profile: the narrative chapters follow the
repository's evidence discipline for historical accounts, and the
philological chapters read each witness at its own locus and preserve
disagreement. It is not a general Mariology, a liturgical book, an
apparition study, or a guide to recovering the object. Excluded or
bounded: contested archaeology is reported at its scholarly confidence
and not adjudicated (Jericho's event archaeology; the Mount Ebal
structure); the \work{Kebra Nagast} and the Aksum claim are treated as
fourteenth-century reception history; modern pseudo-archaeology is
declined in a sentence; Islamic reception is bounded to one paragraph on
Qur'an 2:248; source hypotheses about the ark narratives (Rost and
successors) are noted without adjudication, the narrative standing on the
canonical text; and no mediation title beyond the sense of \work{Lumen
Gentium} 60--62 is advanced.

\textbf{Authority classes, kept distinct.} (1) Scripture supplies the
whole substance: every canonical ark pericope from Exodus 25 to
Apocalypse 11:19 is treated, quoted from the Douay-Rheims base text; the
study adds nothing to public Revelation, and reads the unity of the two
Testaments under \work{Dei Verbum} 12 and CCC 128--130, the rule stated
in the chapter on the rules of typology. (2) Defined dogma touched here: the divine
motherhood (Ephesus, 431) and the Assumption (1950). The Mary-ark
typology itself belongs to no defined formula: \work{Munificentissimus
Deus} deploys the type in its argument (nn.~26, 29, 34) and defines
without ark language (n.~44), and the study claims nothing more.
(3) Conciliar and papal magisterium, cited act by act:
\work{Ineffabilis Deus} (1854); \work{Munificentissimus Deus} with its
\work{Acta} text and notes; \work{Marialis Cultus} 6 and 26; CCC 489,
722, 2676--2677; \work{Lumen Gentium} 60--62 governing everything said of
Mary, with the verified negative that the constitution never uses the
title; John Paul II's 1983 Angelus catechesis beside the silence of
\work{Redemptoris Mater}. (4) Liturgical reception, treated as real
evidence on the warrant stated in the chapter on the rules of typology: dated acts of
two rites, from the Akathist and the Dormition Menaion through the Litany
of Loreto (approved 1587, fixed 1601), the Damascene Matins lesson of the
Roman office (1568--1950, verified here in a 1942 printing), the revised
office of 1950, and the postconciliar Lectionary 621/622 with the
Liturgia Horarum. (5) Patristic witness, carried with per-text
authenticity classes: genuine and dated witnesses (Hippolytus, Ephrem,
Proclus's fifth oration, Hesychius, Chrysippus, Jacob of Serugh, John
Damascene, and Augustine's contrary exegesis); pseudonymous or disputed
texts used only under honest labels (Ps.-Gregory Thaumaturgus,
Ps.-Methodius, the Turin homily attributed to Athanasius, the sermon
transmitted under both Ambrose's and Maximus's names, author unknown);
and verified negatives printed as findings (no genuine Augustine or
Ambrose Marian ark text; no ark in Cyril's Ephesus homily; the silences
of \work{Cogitis me}, Bernard, and Aquinas; the chant-index negative for
the Assumption). (6) Historical scholarship: the ark-origins literature
(de Vaux, Haran, Mettinger, Seow, Eichler, Noegel, and the compositional
school), the excavation records of Shiloh and Jericho, Miller and
Roberts, and the modern Apocalypse commentators (Bauckham, Aune, Beale)
are reported as schools with premises; where a work was consulted at one
remove, the note says so. (7) The Luke debate is presented as genuinely
open and methodological: Burrows, Lyonnet, and Laurentin against Brown,
Fitzmyer, Schürmann, Bovon, and Isaacs, with the Kozłowski and de Koning
revival --- both sides quoted at their strongest, no verdict imposed
beyond what the Visitation section prints. (8) Editorial synthesis,
labeled where it occurs: the study's one framed synthesis box, in the
Assumption section, joins 2 Machabees 2:7--8 to the seventh trumpet; the
identification of Psalm 77:60--61's ``strength'' and ``beauty'' with the
ark is labeled source-grounded synthesis in the prose of the pattern
section; and the master table's prose grades, the fate-traditions table,
and the coda's homiletic extensions are project constructions grounded in
the cited sources and labeled at their places.

\textbf{The two axes.} One promise of method governs the study, made in
the overture and stated formally in the chapter on the rules of typology: facts exact, framework
Catholic, labels honest. Questions of fact are settled by evidence;
questions of interpretation are not settled by any school's minimalism;
liturgical and magisterial reception counts as evidence; and every
disputed reading is declared in one or more of four states ---
philologically anchored, anciently attested, liturgically or
magisterially received, or the project's own labeled synthesis --- under
the named distinction that attestation is not reception, and the absence
of either is not falsity.

\textbf{Edition and as-of bounds.} English Scripture is quoted throughout
from the public-domain Douay-Rheims (Challoner) text as printed at
drbo.org, with its book names (1--4 Kings, 1--2 Paralipomenon, Josue,
Jeremias, 2 Machabees, Apocalypse) and Vulgate psalm numbers, dual
citations given where confusion threatens. Septuagint claims follow
Rahlfs numbering and were verified in a machine-readable derivative of
Rahlfs 1935, independently re-run in an adversarial pass; a spot-check
against the printed Rahlfs--Hanhart remains open below. New Testament
Greek counts use the SBL Greek New Testament. The print contains no
Greek or Hebrew script: Greek is transliterated in italics with macrons,
Hebrew in plain unpointed transliteration; the exact scripts stand in
the research records. All web witnesses were read 2026-08-15 unless a
note states another date (the chant-index data were retrieved in August
2026). Mutable claims --- the current lectionary selections, the
Liturgia Horarum text served by derivative electronic witnesses, the
Triodion rubric reported from modern liturgical resources --- are stated
as of those witnesses and dates and must be rechecked before reliance.

\textbf{Translation status.} Scripture in English is Douay-Rheims only.
Patristic and ancient English quotations use identified public-domain
translations, named at each locus: the Ante-Nicene and
Nicene/Post-Nicene series translators, Morris (1847) for Ephrem, Livius
(1893), Allies (1898) for the Damascene (an identification noted as not
confirmed line by line), Marshall (1878) for Bede, Whiston for Josephus,
Charles for 2 Baruch, Kraft and Purintun for 4 Baruch, Pickthall for the
Qur'an. Quotations from in-copyright scholarly translations (the modern
Apocalypse-commentary volumes; Scheeben's translator) are kept to brief
attributed scale. Texts with no public-domain English --- the Turin
homily, Jacob of Serugh, Gregory the Great's homily, the Latin hymns and
the Breviary lesson, Primasius, the Greek homilists where no
public-domain rendering is credited --- are given in labeled editorial
paraphrase, each marked as a paraphrase and not a translation, or are
described with their loci. Rabbinic, Menaion, and current lectionary
texts are described or paraphrased with exact references; their modern
English translations are not quoted. No prayer or liturgical text is
offered for recitation, and the leaf therefore keeps no prayer-text
audit.

\textbf{Textual cruxes carried.} The base text is quoted as printed and
never silently emended; where the witnesses divide, the division is
printed. The cruxes, each at its section: the Bethsames number, 1 Kings
6:19 (fifty thousand and seventy quoted as printed; the Masoretic
anomaly, Josephus's seventy, the Revised Standard Version's flagged
emendation, and the negative that 4QSam-a does not preserve the verse
--- the catastrophe section); the mice of 1 Kings 5:6 (Vulgate and
Septuagint against the Hebrew --- the same section); the mood and
arithmetic of 1 Kings 7:2 (``rested'' against the Hebrew's lamenting;
the twenty years --- the same section); the word that killed Oza (Douay
``for his rashness'' as the Vulgate's construal of the obscure
\term{shal}, with 4QSam-a and 1 Paralipomenon 13:10 reading the touch
--- the processions section); Nachon against Chidon, and the sacrifices
of the six paces (an ox and a ram against seven bulls, the scroll
siding with Paralipomenon --- the same section, tabled, not
harmonized); the ark's contents (nothing but the tables at 3 Kings 8:9
against the fuller inventory of Hebrews 9:4 --- stated at Sinai,
restated at the temple, untied nowhere); and the pronouns of Psalm
77:60--61 (the Douay's ``their'' against the Hebrew's ``his'' --- the
pattern section, with the ark identification labeled as synthesis).

\textbf{Unresolved and outstanding.} The open verification queue,
disclosed note by note in the body and gathered here: the standing of
Hesychius's Homily V in Aubineau's critical edition of the festal
homilies was not consulted, and the attribution is carried under that
caution; Mutzenbecher's disposition of the doublet sermon remains
unchecked; the Clementine Vulgate's Psalm 131:8 was not collated against
the registered Clementine edition (the same Latin stands verified in the
\work{Acta Apostolicae Sedis} and in Migne's Hesychius column); the
machine-readable Septuagint corpus awaits its spot-check against the
printed Rahlfs--Hanhart; the Visitation responsory \work{Habitavit arca
domini} survives truncated in the chant index, and its Munich manuscript
has not been read; whether the Damascene \term{arca} sentence stood in
the 1568 editio princeps of the breviary lesson was not established; the
a Lapide negative awaits a page-image check of an Antwerp edition; Jacob
of Serugh is cited through Puthuparampil, the wording to be re-verified
against Hansbury before any longer quotation; and the 1558 Dillingen
print of the litany was not inspected. Several modern monographs are
reported at one remove and flagged as such in their notes, including
Laurentin's and Brown's pages as transmitted in Kozłowski, the
ark-origins monographs, Miller and Roberts, and the modern Apocalypse
commentaries. The ark-or-ephod variant at 1 Kings 14:18, noted in the
research scope as a crux to preserve, was in the event not treated in
the body and is recorded here as untreated. Outstanding review:
independent Mariological, patristic, liturgical, historical, and rights
review; collation against print editions; ecclesiastical review or
approval (none claimed). Publication language for this artifact is
``source-audited working reference.''

\textbf{Drawings.} The seven drawings are schematic readings of the
texts they cite, not archaeological reconstructions; the map is not to
scale and takes the traditional site identifications, as its caption
states. Because the drawings cannot yet be carried into the web
rendering path, this study is published in print form only, and its web
edition is withdrawn until that path exists.

\textbf{Records.} Research records for this leaf, under
\texttt{research/}: \texttt{scope.md}; \texttt{source-audit.md};
\texttt{evidence-map.md} (the narrative chapters' evidence map); and
\texttt{source-bindings.toml} (bindings into the repository source
library).
